\documentclass[a4paper,11pt]{article}
\usepackage{exam-style-842-unit}
\examinehead{842 信息化概论 · 第07单元练习册}{练习一\quad 人工智能基础}
\begin{document}

\unitcover{842 信息化概论 · 单元练习册}{第07单元\quad 人工智能技术}{练习一\quad 人工智能基础（共 100 分）}{约 70 分钟}{智能及其本质；人工智能的概念与定义、智能机器、脑智能与群智能、符号智能与计算智能；人工智能发展简史（孕育—形成—知识应用—综合集成）；人工智能主要内容八大类；研究途径与方法；图灵测试；弱人工智能与强人工智能；机器学习、深度学习与神经网络的关系。}

\begin{knowmap}
\textbf{认知链：}数据（记载符号）$\rightarrow$ 信息（减少不确定性）$\rightarrow$ 知识（规律化认识）$\rightarrow$ 智能（获取并运用知识解决问题的能力）；智能的三点共识：感知能力 · 记忆和思维能力 · 学习能力、自适应能力及行为能力\\[2pt]
\textbf{三种定义视角：}尼尔逊（关于知识的学科）· 功能定义（智能机器执行判断、推理、证明、识别、学习、求解等人类智能功能）· 实用定义（知识工程学）；9.2.1 的一般定义：模拟、延伸和扩展人的智能的理论、方法、技术及应用系统\\[2pt]
\textbf{发展四阶段：}孕育（人工神经元模型）$\rightarrow$ 形成（1956 年达特茅斯会议，麦卡锡提议采用“人工智能”术语）$\rightarrow$ 知识应用（1977 年费根鲍姆提出知识工程）$\rightarrow$ 综合集成（20 世纪 80 年代末起）\\[2pt]
\textbf{主要内容八类：}搜索与求解 · 学习与发现 · 知识与推理 · 发明与创造 · 感知与交流 · 记忆与联想 · 系统与建造 · 应用与工程\\[2pt]
\textbf{研究途径：}心理模拟符号推演 · 生理模拟神经计算 · 行为模拟控制进化 · 群体模拟仿生计算 · 博采广鉴自然计算 · 原理分析数学建模\\[2pt]
\textbf{概念关系（教材未展开，按通用教材补充）：}人工神经网络包含于深度学习，深度学习包含于机器学习，机器学习包含于人工智能；弱人工智能只解决特定任务，强人工智能要求通用智能与意识
\end{knowmap}

\begin{keyterms}
\textbf{智能}：人类在认识世界和改造世界的活动中由脑力劳动表现出来的能力，即对客观事物进行合理分析、判断及有目的的行动和有效处理周围环境事宜的综合能力。\\
\textbf{人工智能}：研究、开发用于模拟、延伸和扩展人的智能的理论、方法、技术及应用系统的一门新的技术科学（教材 9.2.1）。\\
\textbf{图灵测试}：图灵 1950 年在《计算机与智能》中提出，让测试者通过对话区分被测试者是人还是机器，若无法区分则认为机器具有智能（教材第 1 章 1.3.4 提及，完整表述按通用教材补充）。\\
\textbf{智能机器}：能够在各类环境中自主地或交互地执行各种拟人任务的机器，至少具备感觉、反应、思考三个要素（教材 9.2.2）。\\
\textbf{脑智能}：宏观心理（语言）层次的智能，是一种个体智能（教材 9.2.3）。\\
\textbf{群智能}：无智能或仅有简单智能的个体通过聚集协同而表现出的智能行为，是一种社会智能或系统智能（教材 9.2.3）。\\
\textbf{符号智能}：模拟脑智能的传统人工智能，以符号形式的知识和信息为基础，主要通过逻辑推理、运用知识进行问题求解（教材 9.2.4）。\\
\textbf{计算智能}：模拟群智能的人工智能，以数值数据为基础，主要通过数值计算、运用算法进行问题求解（教材 9.2.4）。\\
\textbf{知识工程}：1977 年费根鲍姆提出，研究知识获取、知识表示、知识组织与管理和知识运用等技术的学科方向（教材 9.3）。\\
\textbf{搜索与求解}：为了达到某一目标而多次进行某种操作、运算、推理或计算的过程，是人工智能最基本的研究内容（教材 9.4.1）。\\
\textbf{人工神经网络}：由大量人工神经元连接而成的网络，可作为信息和知识的载体，用神经计算实现学习、记忆、联想、识别和推理（教材 9.5.2）。\\
\textbf{机器学习}：让机器从数据与经验中获取知识、发现规律的技术，对应教材“学习与发现”（教材仅简要提及，系统分类按通用教材补充）。\\
\textbf{深度学习}：多层的人工神经网络和训练的方法，通过非线性激活逐层提取特征（教材 9.6.14）。\\
\textbf{弱人工智能与强人工智能}：弱人工智能专注特定任务、不具备真正的理解与意识；强人工智能要求具备通用智能乃至意识（教材未使用该术语，按通用教材补充）。
\end{keyterms}

\begin{contrast}{脑智能 / 群智能（每行依次为：核心含义·层次定位·典型例子）}
\citem{脑智能}{宏观心理（语言）层次上的智能}{个体智能}{人的记忆、联想、推理、计算、思考等心理活动}
\citem{群智能}{无智能或简单智能的个体经聚集协同表现出的智能行为}{社会智能或系统智能}{蚁群觅食、鸟群飞翔；神经元群活动形成脑智能}
\end{contrast}

\begin{contrast}{符号智能 / 计算智能（每行依次为：模拟对象·基础与求解方式·主要内容）}
\citem{符号智能}{模拟脑智能（传统人工智能、经典人工智能）}{以符号形式的知识和信息为基础，主要经逻辑推理、运用知识求解}{知识获取、知识表示、知识组织与管理、知识运用}
\citem{计算智能}{模拟群智能}{以数值数据为基础，主要经数值计算、运用算法求解}{神经计算、进化计算、免疫计算、蚁群算法、粒群算法、自然计算、人工生命}
\end{contrast}

\begin{contrast}{符号主义 / 连接主义 / 行为主义（每行依次为：模拟层面·核心方法·代表成果）}
\citem{心理模拟，符号推演}{人脑的宏观心理层面}{把问题或知识表示成逻辑网络，采用符号推演模拟逻辑思维}{基于知识的智能系统、专家系统}
\citem{生理模拟，神经计算}{人脑的生理层面（微观结构和工作机理）}{用人工神经网络作载体，用神经计算实现学习、记忆、联想、识别与推理}{人工神经网络}
\citem{行为模拟，控制进化}{人与动物在和环境交互、控制中的行为}{基于“感知—行为”模型模拟反应、适应、学习、寻优}{布鲁克斯的六足行走机器人}
\end{contrast}

\begin{contrast}{机器学习 / 深度学习 / 人工神经网络（每行依次为：概念定位·相互关系·典型代表）}
\citem{机器学习}{让人工智能系统从数据与经验中获取知识、发现规律}{范围最广，包含深度学习等多种方法}{决策树、支持向量机、神经网络等}
\citem{深度学习}{多层人工神经网络及其训练方法}{属于机器学习，以多层神经网络为模型}{AlphaGo 的两个神经网络“大脑”}
\citem{人工神经网络}{由大量人工神经元连接而成的计算模型}{是深度学习的模型基础，但神经网络不等同于深度学习}{多层感知机、卷积神经网络等}
\end{contrast}

\qtypename{一、单项选择题}{10 题，每题 3 分，共 30 分}
\q 一般认为，智能是（\quad）
\choices{A. 人类在认识世界和改造世界的活动中由脑力劳动表现出来的能力}{B. 计算机自动执行程序的能力}
{C. 生物对外界刺激的本能反应}{D. 一个人掌握的知识总量}

\q 提出“人工智能是关于知识的学科，是怎样表示知识以及怎样获得知识并使用知识的科学”这一观点的是（\quad）
\choicesline{A. 图灵}{B. 尼尔逊}{C. 麦卡锡}{D. 费根鲍姆}

\q 图灵测试的基本做法是（\quad）
\choices{A. 测试机器的运算速度是否超过人脑}{B. 让测试者通过对话区分被测试者是人还是机器，若无法区分则认为机器具有智能}
{C. 检查机器能否完成数学定理证明}{D. 测量机器的存储容量和通信带宽}

\q “人工智能”这一术语正式提出于（\quad）
\choices{A. 20 世纪 40 年代提出人工神经元模型时}{B. 1950 年图灵发表《计算机与智能》时}
{C. 1956 年达特茅斯学院会议}{D. 1977 年第五届国际人工智能联合会议}

\q 专家系统大量涌现、商品化智能系统纷纷推出，属于人工智能发展简史中的（\quad）
\choicesline{A. 孕育期}{B. 形成期}{C. 知识应用期}{D. 综合集成期}

\q 1977 年在第五届国际人工智能联合会议上正式提出知识工程概念的是（\quad）
\choicesline{A. 麦卡锡}{B. 费根鲍姆}{C. 布鲁克斯}{D. 尼尔逊}

\q 教材指出，人工智能最基本的研究内容是（\quad）
\choicesline{A. 搜索与求解}{B. 发明与创造}{C. 系统与建造}{D. 应用与工程}

\q 下列各项中，属于计算智能主要内容的是（\quad）
\choicesline{A. 知识表示}{B. 谓词逻辑推理}{C. 蚁群算法}{D. 专家系统}

\q 按教材 9.2.2，智能机器至少要具备的三个要素是（\quad）
\choices{A. 感觉要素、反应要素、思考要素}{B. 感知要素、记忆要素、学习要素}
{C. 输入要素、处理要素、输出要素}{D. 动力要素、传动要素、控制要素}

\q 下列关于人工智能、机器学习、深度学习、神经网络关系的表述，正确的是（\quad）
\choices{A. 人工智能就是机器学习}{B. 深度学习就是神经网络，二者完全等同}
{C. 神经网络包含机器学习}{D. 深度学习是机器学习的一个分支，通常以多层人工神经网络为模型}

\qtypename{二、多项选择题}{4 题，每题 4 分，共 16 分；每题有两个或两个以上正确项}
\q 教材所述智能本质的共识包括（\quad）
\choicesline{A. 感知能力}{B. 记忆和思维能力}{C. 学习、自适应及行为能力}{D. 情感能力}

\q 符号智能的主要内容包括（\quad）
\choicesline{A. 知识获取}{B. 知识表示}{C. 知识组织与管理}{D. 数值计算与优化搜索}

\q 关于人工智能的研究途径与方法，下列说法正确的有（\quad）
\choices{A. 心理模拟符号推演从人脑的宏观心理层面入手}{B. 生理模拟神经计算以人工神经网络作为信息和知识的载体}
{C. 行为模拟法的典型代表是布鲁克斯研制的六足行走机器人}{D. 群体模拟仿生计算只研究单个神经元的工作机理}

\q 关于人工智能的发展简史，下列说法正确的有（\quad）
\choices{A. 1956 年达特茅斯学院会议正式确立了人工智能的研究领域}{B. 1977 年费根鲍姆正式提出知识工程的概念}
{C. 1969 年国际人工智能联合会议标志着人工智能得到国际认可}{D. 人工智能自诞生后发展一直顺利，从未出现高潮与低谷}

\qtypename{三、判断题}{6 题，每题 2 分，共 12 分；正确的画“√”，错误的画“×”并写出改正后的表述}
\q 智能是先天遗传的能力，与后天的学习和实践无关。\quad（\quad）\quad 改正：\fixblank

\q 脑智能是个体智能，群智能是社会智能（系统智能）。\quad（\quad）\quad 改正：\fixblank

\q 符号智能以数值数据为基础，主要通过数值计算、运用算法进行问题求解。\quad（\quad）\quad 改正：\fixblank

\q 人工智能学科正式诞生于 1956 年的达特茅斯学院会议。\quad（\quad）\quad 改正：\fixblank

\q 图灵测试通过就证明机器具有与人相同的意识和情感。\quad（\quad）\quad 改正：\fixblank

\q 深度学习是多层的人工神经网络及训练方法，可视为机器学习的一个分支。\quad（\quad）\quad 改正：\fixblank

\qtypename{四、配对题}{6 对，每对 1 分，共 6 分；把 A 栏与 B 栏对应内容配对}
\q 将下列人物、算法或概念与其对应内容配对，把 B 栏字母填入括号。
\par\vspace{0.2em}
\noindent\begin{tabular}{@{}p{0.44\textwidth}@{\hspace{1.2em}}p{0.44\textwidth}@{}}
（1）麦卡洛克与匹茨 \hfill（\quad） & A. 1977 年提出知识工程 \\
（2）麦卡锡 \hfill（\quad） & B. 研制六足行走机器人，行为模拟的代表 \\
（3）费根鲍姆 \hfill（\quad） & C. 提出人工神经元模型，人工智能孕育期的最早工作 \\
（4）布鲁克斯 \hfill（\quad） & D. 模拟蚂蚁群体觅食行为的仿生计算 \\
（5）蚁群算法 \hfill（\quad） & E. 1956 年提议正式采用“人工智能”这一术语 \\
（6）符号智能 \hfill（\quad） & F. 以符号知识为基础、主要经逻辑推理求解 \\
\end{tabular}\par\vspace{0.3em}

\qtypename{五、简答题}{3 题，每题 8 分，共 24 分}
\q 简述智能的本质（教材概括的三点共识），并指出思维分为哪三种。\anslines{2.4cm}

\q 人工智能有哪几种定义视角？教材 9.2.1 给出的一般定义是什么？智能机器至少应具备哪三个要素？\anslines{2.4cm}

\q 简述教材 9.5 所述人工智能的研究途径与方法（写出五种并各举一例）。\anslines{2.4cm}

\newpage
\qtypename{六、综合题}{1 题，共 12 分}
\q 某农业科技公司开发了一套作物病虫害识别系统：用大量标注过的叶片图像训练多层神经网络，用户拍照上传即可得到病虫害名称与防治建议。公司宣传该系统是“具有人类智能的强人工智能”。请结合教材回答：（1）该系统以数值数据为基础、通过算法训练求解，属于符号智能还是计算智能？说明理由。（2）该系统主要涉及人工智能八大主要内容中的哪几类？（3）说明机器学习、深度学习与人工神经网络三者的关系，并指出该系统属于哪一类。（4）从强人工智能的判断标准看，公司的宣传是否成立？说明理由。（5）结合人工智能发展简史，说明该系统处于怎样的技术背景。\anslines{5cm}

\end{document}
